\section*{有限维模}

\begin{enumerate}
\def\labelenumi{\arabic{enumi})}
\setcounter{enumi}{18}
\item
  令 V 为有限维 g-模。对所有 \(\mu\in h^{*}\)，令 \(V^{\mu}\) 为由 \(v\in V\)（满足 \(h.v=\mu(h)v\) 对所有 \(h\in h\)）组成的集合。\(V^{\mu}\) 的维数称为 \(\mu\) 在 V 中的重数；若它 \(\geq1\)，即若 \(V^{\mu}\neq0\)，则称 \(\mu\) 是 V 的权。我们有 \(V=\bigoplus_{\mu\in h^{*}}V^{\mu}\)。V 的每个权都属于 P(R)。若 \(\mu\) 是 V 的权且 \(w\in W\)，则 \(w\mu\) 也是 V 的权，且重数与 \(\mu\) 相同。若 \(v\in V^{\mu}\) 且 \(x\in g^{\alpha}\)，则 \(x.v\in V^{\mu+\alpha}\)。
\item
  令 B 为 R 的基。给定 B 就确定了 \(h_{Q}^{*}\) 上的一个序关系：\(h_{Q}^{*}\) 中 \(\geq 0\) 的元素是 B 中元素以 \(\geq 0\) 的有理数为系数的线性组合。以 \(Q_{+}(R)\)（分别地，\(R_{+}\)）表示 \(Q(R)\)（分别地，R）的正元素所成的集合。
\end{enumerate}

令 \(V\) 为有限维单 \(\mathfrak{g}\)-模。则 \(V\) 有最高权 \(\lambda\)。这个权的重数为 1，且它是支配权：若 \(\alpha \in R_{+}\)，\(\lambda(H_{\alpha})\) 是 \(\geq 0\) 的整数。我们有 \(\mathfrak{g}^{\alpha}V^{\lambda} = 0\)（若 \(\alpha \in R_{+}\)）。\(V\) 的每个权都具有 \(\lambda - \nu\)（\(\nu \in Q_{+}(R)\)）的形式；反之，若一个权是支配权且具有 \(\lambda - \nu\)（\(\nu \in Q_{+}(R)\)）的形式，则它是 \(V\) 的权。

\begin{enumerate}
\def\labelenumi{\arabic{enumi})}
\setcounter{enumi}{20}
\tightlist
\item
  具有相同最高权的两个有限维单 g-模同构。每个支配权都是某个有限维单 g-模的最高权。
\end{enumerate}

每个有限维单 g-模都是绝对单的。

\begin{enumerate}
\def\labelenumi{\arabic{enumi})}
\setcounter{enumi}{21}
\item
  令 \(\varPhi\) 为 g 的 Killing 型，\(C \in U(\mathfrak{g})\) 为相应的 Casimir 元，\(\langle\cdot,\cdot\rangle\) 为 \(h^{*}\) 上的、\(\varPhi|h \times h\) 的逆形式，\(\rho = \frac{1}{2} \sum_{\alpha \in R_{+}} \alpha\)。令 V 为有限维单 g-模，最高权为 \(\lambda\)。则 \(C_{V}\) 是比为 \(\langle\lambda, \lambda + 2\rho\rangle\) 的位似。
\item
  令 \(\mathrm{V}\) 为有限维 \(\mathfrak{g}\)-模，\(\mathrm{V}^*\) 为其对偶。则 \(\mu \in \mathfrak{h}^*\) 是 \(\mathrm{V}^*\) 的权当且仅当 \(-\mu\) 是 \(\mathrm{V}\) 的权，且 \(\mu \in \mathrm{V}^*\) 的重数等于 \(-\mu\) 在 V 中的重数。若 V 是最高权为 \(\lambda\) 的单模，则 \(V^{*}\) 是最高权为 \(-w_{0}\lambda\) 的单模，其中 \(w_{0}\) 是 W 中把 B 变为 -B 的元素。
\item
  令 \(\mathrm{V}\) 为有限维单 \(\mathfrak{g}\)-模，其最高权为 \(\lambda\)，\(\mathcal{B}\) 为 \(\mathfrak{g}\)-不变双线性形式（\(\mathrm{V}\) 上的）所成的向量空间。令 \(m\) 为整数 \(\sum_{\alpha \in \mathbb{R}_+} \lambda(H_\alpha)\)，并令 \(w_0 \in \mathrm{W}\) 如 23) 中那样。若 \(w_0 \lambda \neq -\lambda\)，则 \(\mathrm{V}\) 与 \(V^*\) 不同构，且 \(\mathcal{B} = 0\)。若 \(w_0 \lambda = -\lambda\)，则 \(\dim \mathcal{B} = 1\)，且 \(\mathcal{B}\) 的每个非零元素都非退化；若 \(m\) 为偶（分别地，奇）数，则 \(\mathcal{B}\) 的每个元素都是对称（分别地，交错）的。
\item
  令 \(\mathbf{Z}[\mathrm{P}]\) 为群 \(\mathrm{P} = \mathrm{P}(\mathrm{R})\) 的系数在 \(\mathbf{Z}\) 中的群代数。若 \(\lambda \in \mathrm{P}\)，以 \(e^{\lambda}\) 表示 \(\mathbf{Z}[\mathrm{P}]\) 中相应的元素；诸 \(e^{\lambda}\)（\(\lambda \in \mathrm{P}\)）构成一个 \(\mathbf{Z}\)-基（\(\mathbf{Z}[\mathrm{P}]\) 的），且 \(e^{\lambda}e^{\mu} = e^{\lambda +\mu}\)（对 \(\lambda, \mu \in \mathrm{P}\)）。
\end{enumerate}

令 \(\mathrm{V}\) 为有限维 \(\mathfrak{g}\)-模。\(\mathrm{V}\) 的特征标记作 \(\operatorname{ch} \mathrm{V}\)，它是元素 \(\sum_{\mu \in \mathrm{P}} (\dim \mathrm{V}^{\mu}) e^{\mu}\)（\(\mathbf{Z}[\mathrm{P}]\) 的）；这个元素属于子代数 \(\mathbf{Z}[\mathrm{P}]^{\mathrm{W}}\)——\(\mathbf{Z}[\mathrm{P}]\) 中由在 \(\mathrm{W}\) 下不变的元素组成的。我们有

\[
\operatorname{ch} \left(\mathrm{V} \oplus \mathrm{V} ^ {\prime}\right) = \operatorname{ch} \mathrm{V} + \operatorname{ch} \mathrm{V} ^ {\prime} \quad \text { and } \quad \operatorname{ch} \left(\mathrm{V} \otimes \mathrm{V} ^ {\prime}\right) = (\operatorname{ch} \mathrm{V}). (\operatorname{ch} \mathrm{V} ^ {\prime}).
\]

具有相同特征标的两个有限维 g-模同构。

对所有 \(\alpha\in B\)，令 \(V_{\alpha}\) 为以基本权 \(\varpi_{\alpha}\)（对应于 \(\alpha\)）为最高权的单 g-模。诸元素 ch \(V_{\alpha},\alpha\in B\) 代数无关，且生成 Z-代数 \(Z[P]^{W}\)。

\begin{enumerate}
\def\labelenumi{\arabic{enumi})}
\setcounter{enumi}{25}
\tightlist
\item
  令 \(\rho\) 为诸根 \(\geq 0\) 的根之和的一半。对所有 \(w \in W\)，令 \(\varepsilon(w)\) 为 w 的行列式，等于 \(\pm 1\)。若 V 是最高权为 \(\lambda\) 的有限维单 g-模，则
\end{enumerate}

\[
\left(\sum_ {w \in \mathrm{W}} \varepsilon (w) e ^ {w \rho}\right). \mathrm{chV} = \sum_ {w \in \mathrm{W}} \varepsilon (w) e ^ {w (\lambda + \rho)},
\]

以及

\[
\dim V = \prod_ {\alpha \in R _ {+}} \frac {\langle \lambda + \rho , H _ {\alpha} \rangle}{\langle \rho , H _ {\alpha} \rangle}.
\]

\begin{enumerate}
\def\labelenumi{\arabic{enumi})}
\setcounter{enumi}{26}
\tightlist
\item
  对所有 \(\nu \in \mathrm{P}\)，令 \(\mathfrak{P}(\nu)\) 为族 \((n_{\alpha})_{\alpha \in \mathrm{R}_{+}}\) 的个数，其中诸 \(n_{\alpha}\) 是 \(\geq 0\) 的整数且满足 \(\nu = \sum_{\alpha \in \mathrm{R}_{+}} n_{\alpha} \alpha\)。令 \(\mathrm{V}\) 为有限维单 \(\mathfrak{g}\)-模，其最高权为 \(\lambda\)。若 \(\mu \in \mathrm{P}\)，则 \(\mu\) 在 \(\mathrm{V}\) 中的重数是
\end{enumerate}

\[
\sum_ {w \in \mathrm{W}} \varepsilon (w) \mathfrak {P} (w (\lambda + \rho) - (\mu + \rho)).
\]

\begin{enumerate}
\def\labelenumi{\arabic{enumi})}
\setcounter{enumi}{27}
\tightlist
\item
  令 \(\mathrm{V}, \mathrm{V}', \mathrm{V}''\) 为有限维单 \(\mathfrak{g}\)-模，\(\lambda, \mu, \nu\) 为它们的最高权。在 \(\mathrm{V} \otimes \mathrm{V}'\) 中，\(\mathrm{V}''\) 型的同构型分支的长度为
\end{enumerate}

\[
\sum_ {w, w ^ {\prime} \in \mathrm{W}} \varepsilon (w w ^ {\prime}) \mathfrak {P} (w (\lambda + \rho) + w ^ {\prime} (\mu + \rho) - (\nu + 2 \rho)).
\]

特别地，若 \(\nu = \lambda + \mu\)，则所讨论的同构型分支是单的，且由 \((\mathrm{V} \otimes \mathrm{V}')^{\lambda + \mu} = \mathrm{V}^{\lambda} \otimes \mathrm{V}'^{\mu}\) 生成。

